\documentclass[10pt]{article}

\def\Type{LessonCaelan}
\def\TypeNum{Lesson 19}
\def\TypeTitle{Summations}
\def\Semester{Fall 2021} %Not Used 
\def\Course{MATH 1190 \& 1210}
\def\Targets{  \bf\color{Fuchsia}I3, \bf\color{Fuchsia}I4 }

\usepackage{preambleDJ}

%%%%%%%%%%%% BEGIN DOCUMENT %%%%%%%%%%%%%%%
%%%%%%%%%%%% BEGIN DOCUMENT %%%%%%%%%%%%%%%
%%%%%%%%%%%% BEGIN DOCUMENT %%%%%%%%%%%%%%%
%%%%%%%%%%%% BEGIN DOCUMENT %%%%%%%%%%%%%%%
%%%%%%%%%%%% BEGIN DOCUMENT %%%%%%%%%%%%%%%

%%%%%%%%%%%% BEGIN DOCUMENT %%%%%%%%%%%%%%%
%%%%%%%%%%%% BEGIN DOCUMENT %%%%%%%%%%%%%%%
%%%%%%%%%%%% BEGIN DOCUMENT %%%%%%%%%%%%%%%
%%%%%%%%%%%% BEGIN DOCUMENT %%%%%%%%%%%%%%%
%%%%%%%%%%%% BEGIN DOCUMENT %%%%%%%%%%%%%%%

\begin{document}
\pagestyle{otherpages}
\thispagestyle{firstpage}
\vs{.3}

\textbf{Intended Learning Outcomes}
After this lesson, you will be able to 

\begin{itemize}
    \item Compute expressions involving summation notation. (\target{F3})
    \item Write sums concisely using summation notation. (\target{F3})
    \item Understand and apply the linearity properties of summation (\target{F3})
    %\item Use summation properties to expand and simplify sums (\target{F3})
\end{itemize}

\vspace{-0.6cm}

\hrulefill\\

%Some important properties of probability rely on being able to add several, similarly-formatted quantities together. For example, one important property of probability is that when we sum together the probabilities of each outcome in the sample space, we should get 1 since there is a 100\% chance that {\it something} in the sample space will happen:
%\[ \text{If }S=\{outcome_1,outcome_2,\ldots,outcome_n\},\text{ then }P(\{outcome_1\}) + P(\{outcome_2\}) + \ldots + P(\{outcome_n\}) = 1.\]
%As we'll see in the coming lessons, the above statement can be more succinctly written as a summation:
%\[ \sum_{i=1}^{n} P(\{outcome_n\}) = 1 \quad\text{ or }\quad \sum_{outcome\in S} P(outcome) = 1\]
In this class, we'll practice using notation for expressing sums using summation so that we will be able to use it in our upcoming probability lessons.

{\example / \exercise}\ltrm{\bf\color{ProcessBlue}F3} Expand and compute the following sums:
    \begin{enumerate}
    \item[(a)] ({\bf\color{blue}Example}) $\displaystyle\sum_{i=0}^2 \frac{2i+1}{i+1}$
        \vfill
    \item[(b)] ({\bf\color{blue}Exercise}) $\displaystyle\sum_{j=1}^4 \log_2\left( 2j\right)$
        \vfill
    \item[(c)] ({\bf\color{blue}Exercise}) $\displaystyle\sum_{k=0}^3 k(-1)^k$
        \vfill
    \end{enumerate}
    
{\example / \exercise}\ltrm{\bf\color{ProcessBlue}F3} Write - but don't compute - the following sums using summation notation:
    \begin{enumerate}
    \item[(a)] ({\bf\color{blue}Example}) $2+3+4+5+6$
        \vfill
    \item[(b)] ({\bf\color{blue}Exercise}) $\displaystyle\frac12+\frac14+\frac16+\frac18+\frac1{10}+\frac1{12}$
        \vfill
    \item[(c)] ({\bf\color{blue}Exercise}) $\displaystyle-\frac15+\frac2{25}-\frac3{125}+\frac4{625}$
        \vfill
    \end{enumerate}
    
\newpage

\important{Summation Properties}
{
Let $c$ be any constant, and let $m$ and $n$ be whole numbers with $m<n$.
\begin{itemize}
\item $\displaystyle\sum_{i=m}^n(a_i+b_i) = \sum_{i=m}^n a_i + \sum_{i=m}^n b_i$ \hfill (linearity property \#1)
\item $\displaystyle\sum_{i=m}^n ca_i = c\sum_{i=m}^n a_i$  \hfill (linearity property \#2)
\item $\displaystyle\sum_{i=1}^n c = cn$
\end{itemize}
}

{\example}\ltrm{\bf\color{ProcessBlue}F3} Use the properties of summations exhaustively to decompose the following summation: \[\sum_{j=1}^n (a_j-3b_j+7c_j-4)\]

    \vfill
    
{\exercise}\ltrm{\bf\color{ProcessBlue}F3} Use the properties of summations exhaustively to decompose the following summations:
    \begin{enumerate}
    \item[(a)] $\displaystyle\sum_{i=1}^n (2a_i+3b_i-1)$
        \vfill
    \item[(b)] $\displaystyle\sum_{k=1}^n \left(\frac{a_k\cdot k}4+9b_k^2-2c_k\cdot d_k\right)$\quad(\emph{Hint:} be careful - this is a bit of a trap.)
        \vfill
    \end{enumerate}

\end{document}